\section{Results — Governance and Oversight}

\subsection{England (Office for Students and Governance System)}

England’s regulator, the Office for Students (OfS), operates a registration-based framework
under the Higher Education and Research Act (HERA). Providers registered with the OfS must
comply with Conditions of Registration A--G, which cover access and participation, quality
and standards, financial sustainability, transparency, information requirements, and student
protection arrangements. 

The registered (lead) provider remains responsible for quality and standards wherever
provision is delivered, including validated, franchised, subcontracted, or transnational
education (TNE) arrangements. Sanctions for non-compliance range from specific conditions to
suspension or deregistration. As an arm’s-length body sponsored by the Department for
Education (DfE), the OfS is governed by a framework document that outlines its strategic
functions, accountability routes, and regulatory powers.

\begin{tcolorbox}
\textbf{Jurisdiction:} England \\
\textbf{Regulator:} Office for Students (OfS) \\
\textbf{Model:} Registration system (Conditions A--G) \\
\textbf{Partnership Responsibility:} Lead provider fully accountable \\
\textbf{Student Protection:} Student Protection Plans (SPP); consumer compliance \\
\textbf{Sanctions:} Targeted conditions; fines; suspension; deregistration \\
\textbf{Key DocIDs:} OFS-22-001; DFE-23-001
\end{tcolorbox}

\subsection{Devolved Nations (Scotland, Wales, Northern Ireland)}

Oversight in the devolved nations differs structurally but retains alignment with UK-wide
expectations for academic standards.

\subsubsection*{Scotland (Scottish Funding Council)}

In Scotland, the Scottish Funding Council (SFC) operates an enhancement-led quality model for
SFC-funded institutions. Rather than a registration system, SFC guidance specifies
quality-enhancement principles, roles, annual evaluation mechanisms, and expectations for
public information. Institutions remain responsible for assuring collaborative or
partnership provision. SFC interventions are primarily linked to funding or support measures.

\begin{tcolorbox}
\textbf{Jurisdiction:} Scotland \\
\textbf{Regulator/Body:} Scottish Funding Council (SFC) \\
\textbf{Model:} Enhancement-led quality approach \\
\textbf{Provider Register:} No \\
\textbf{Partnership Responsibility:} Institutions assure collaborative provision \\
\textbf{Student Interests:} Student engagement; public information \\
\textbf{Sanctions/Interventions:} Funding conditions; support interventions \\
\textbf{Key DocID:} SFC-25-001
\end{tcolorbox}

\subsubsection*{Wales (CTER / Welsh Government)}

Wales operates under the Commission for Tertiary Education and Research (CTER), which sets
policy for registration, designation for student support, and quality assessment frequency.
Providers remain responsible for partnership delivery, and CTER may issue regulatory
directions or amend funding conditions.

\begin{tcolorbox}
\textbf{Jurisdiction:} Wales \\
\textbf{Regulator/Body:} CTER / Welsh Government \\
\textbf{Model:} Registration and designation policy \\
\textbf{Provider Register:} Yes \\
\textbf{Partnership Responsibility:} Provider accountable \\
\textbf{Quality Assessment:} Frequency signalled by CTER \\
\textbf{Sanctions/Interventions:} Directions; variation/withdrawal of funding \\
\textbf{Key DocID:} WELSH-25-001
\end{tcolorbox}

\subsubsection*{Northern Ireland (Department for the Economy)}

Northern Ireland does not operate a dedicated registration system; oversight is provided
through DfE~NI policy frameworks linked to funding, accountability, transparency, and
strategic objectives. Institutions remain responsible for the quality and standards of any
collaborative provision under their awarding powers.

\begin{tcolorbox}
\textbf{Jurisdiction:} Northern Ireland \\
\textbf{Regulator/Body:} Department for the Economy (DfE~NI) \\
\textbf{Model:} Policy and funding oversight \\
\textbf{Provider Register:} No \\
\textbf{Partnership Responsibility:} Institutions accountable for provision \\
\textbf{Student Interests:} Transparency and public information \\
\textbf{Interventions:} Funding or designation changes \\
\textbf{Key DocID:} NIDE-25-001
\end{tcolorbox}

\subsection{Synthesis Across UK Nations}

Across the UK, governance varies in regulatory structure but remains aligned in protecting
academic standards, student interests, and institutional accountability:

\begin{itemize}
    \item England’s OfS model is the most prescriptive, using conditions of registration.  
    \item Scotland emphasises enhancement and funding-linked quality arrangements.  
    \item Wales integrates oversight into a unified tertiary framework via CTER.  
    \item Northern Ireland sets expectations through policy, funding, and accountability routes.  
\end{itemize}

Despite these differences, all nations rely on sector-recognised standards, programme-level
academic regulations, and institutional responsibility for partnership delivery. This
creates a coherent set of expectations for UK providers delivering transnationally,
including in Sri Lanka.
